\section{Graded Exercise Bank}
\label{sec:ch2-exercise-bank}

The exercises are graded by mathematical maturity rather than by algebraic
length.  Basic problems establish fluency; intermediate problems combine
ideas; advanced problems emphasize proof and structure; challenge problems
connect the chapter to mechanics, field theory, and quantum mechanics.
\subsection{Basic fluency}
\begin{enumerate}
  \item State the domain and range of $f(x)=\sqrt{9-x^2}$.
  \item Sketch the level sets of $f(x,y)=x+y$ for levels $-1,0,1$.
  \item Evaluate $\lim_{x\to3}(x^2-9)/(x-3)$.
  \item Evaluate $\lim_{h\to0}[(2+h)^3-8]/h$.
  \item Determine whether $f(x)=|x|$ is differentiable at zero.
  \item Differentiate $x^4-3x^2+7$.
  \item Differentiate $\sin(x^2)$.
  \item Differentiate $xe^x$.
  \item Find the tangent line to $y=\ln x$ at $x=1$.
  \item For $x(t)=5t^2-2t$, find velocity and acceleration.
  \item Compute $f_x$ and $f_y$ for $f=x^3y^2$.
  \item Compute all first partial derivatives of $f=e^{xyz}$.
  \item Find $df$ for $f=x^2+y^2$.
  \item Linearize $f(x,y)=\sqrt{x+y}$ at $(2,2)$.
  \item Find the tangent plane to $z=x^2+y^2$ at $(1,1,2)$.
  \item Compute $\nabla f$ for $f=xyz$.
  \item Find $D_{\hat u}f$ for $f=x^2+y^2$ at $(1,0)$ with $\hat u=(0,1)$.
  \item Find the direction of steepest increase of $f=3x-y$.
  \item Compute the Jacobian of $F(x,y)=(x+y,x-y)$.
  \item Compute the determinant of the Jacobian in Exercise 2.19.
  \item Use the chain rule for $z=x^2+y^2$, $x=\cos t$, $y=\sin t$.
  \item Find the critical points of $f=x^2+y^2$.
  \item Classify the origin for $f=x^2-y^2$.
  \item Compute $\nabla\cdot(x,y,z)$.
  \item Compute $\nabla\times(-y,x,0)$.
  \item Compute $\nabla^2(x^2+y^2+z^2)$.
  \item Differentiate $r=\sqrt{x^2+y^2+z^2}$ with respect to $x$.
  \item Compute $\nabla r$ for $r\ne0$.
  \item Show that the directional derivative is linear in the direction vector.
  \item Explain why a non-unit direction vector changes the numerical value of a directional derivative.
\end{enumerate}
\subsection{Intermediate synthesis}
\begin{enumerate}
  \item Use an $\varepsilon$--$\delta$ argument to prove $\lim_{x\to2}(3x-1)=5$.
  \item Investigate the limit of $x^2y/(x^4+y^2)$ at the origin along several paths.
  \item Show that $f(x,y)=xy/\sqrt{x^2+y^2}$ is continuous at the origin when defined there as zero.
  \item Derive the quotient rule from the product and chain rules.
  \item Find $d^2y/dx^2$ for the implicit curve $x^2+y^2=1$.
  \item Estimate $1/1.98$ using linearization at $x=2$.
  \item Find the second-order Taylor polynomial of $e^x$ at zero.
  \item Compute the Hessian of $f=x^2y+xy^2$.
  \item Find all tangent directions to the level curve $x^2+xy+y^2=3$ at $(1,1)$.
  \item For $T=e^{-x^2-y^2}$, find the direction of fastest cooling at $(1,1)$.
  \item Find the normal line to $x^2+y^2+z^2=9$ at $(1,2,2)$.
  \item Compute the Jacobian of spherical coordinates.
  \item Show that the polar Jacobian vanishes at the coordinate singularity $r=0$.
  \item For $F(x,y)=(e^x\cos y,e^x\sin y)$, compute and interpret $\det J_F$.
  \item Use the chain rule to differentiate $f(r(t))$ for $f=x^2y$ and $r(t)=(t,e^t)$.
  \item Classify all critical points of $f=x^4+y^4-4xy$.
  \item Find extrema of $x+y$ on the circle $x^2+y^2=1$.
  \item Find the point on $z=x^2+y^2$ nearest $(0,0,3)$.
  \item Compute divergence and curl of $F=(yz,zx,xy)$.
  \item Verify directly that $\nabla\times(\nabla f)=0$ for $f=x^2yz$.
  \item Verify directly that $\nabla\cdot(\nabla\times F)=0$ for $F=(xy,yz,zx)$.
  \item Compute $\nabla^2(r^2)$ in three dimensions.
  \item Compute $\nabla^2(1/r)$ for $r\ne0$.
  \item Derive first-order relative-error propagation for $Q=ab^2/c$.
  \item Estimate the change in kinetic energy $K=mv^2/2$ under small changes $dm,dv$.
\end{enumerate}
\subsection{Advanced analysis}
\begin{enumerate}
  \item Prove that differentiability at a point implies continuity there.
  \item Give a function whose partial derivatives exist at the origin but which is not continuous there.
  \item Prove the multivariable chain rule using remainder terms.
  \item Show that the derivative, when it exists, is unique.
  \item Derive the Hessian transformation rule under an orthogonal change of coordinates.
  \item Prove that $\|\nabla f\|$ is the maximal unit directional derivative.
  \item Derive the second-derivative test from the quadratic Taylor approximation.
  \item Find and classify constrained extrema of $xyz$ on $x^2+y^2+z^2=3$.
  \item Derive the Jacobian for cylindrical coordinates and interpret each scale factor.
  \item Show that a continuously differentiable map with nonzero Jacobian determinant is locally orientation preserving or reversing according to its sign.
  \item For a radial field $F=f(r)\hat r$, derive $\nabla\cdot F$ for $r\ne0$.
  \item For a radial scalar $u(r)$, derive $\nabla^2u=u''+2u'/r$ in three dimensions.
  \item Show that $e^{i\mathbf k\cdot\mathbf r}$ is an eigenfunction of the Laplacian.
  \item Derive the material derivative $D/Dt=\partial_t+\mathbf v\cdot\nabla$.
  \item Use dimensional analysis to check every term in the diffusion equation $\partial_tT=\kappa\nabla^2T$.
\end{enumerate}
\subsection{Challenge and physics bridge}
\begin{enumerate}
  \item Construct a function continuous at the origin whose directional derivatives all exist there but which is not differentiable there.
  \item Prove the inverse-function theorem in one dimension and explain which step fails when $f'(a)=0$.
  \item Derive the Euler--Lagrange equation from a first variation of the action.
  \item Show that Hamilton's equations imply conservation of $H$ when $\partial H/\partial t=0$.
  \item Derive Liouville's theorem by computing the divergence of Hamiltonian phase-space flow.
  \item Use the Hessian to derive the normal-mode approximation near a stable equilibrium.
  \item Derive the probability current for a Schrödinger wavefunction and identify where gradients enter.
  \item Show how gauge transformation $A\mapsto A+\nabla\chi$ leaves $B=\nabla\times A$ unchanged.
  \item Explain why $\nabla^2(1/r)=0$ away from the origin yet represents a point source distributionally.
  \item Derive the local conservation equation from a balance law on an arbitrary small volume.
\end{enumerate}

\subsection{Selected hints and solutions}
\begin{description}
\item[Basic 14.] Use $df=f_x\,dx+f_y\,dy$ at $(2,2)$; the denominator in both partial derivatives is $2\sqrt{x+y}$.
\item[Basic 19.] The matrix is $\begin{pmatrix}1&1\\1&-1\end{pmatrix}$; its determinant is $-2$, so the map reverses orientation and doubles area locally.
\item[Intermediate 2.] Along $y=mx^2$ the expression becomes $m/(1+m^2)$, which depends on $m$; therefore the limit does not exist.
\item[Intermediate 10.] $\nabla T=-2e^{-x^2-y^2}(x,y)$; fastest cooling is in the direction $-\nabla T$.
\item[Intermediate 12.] The absolute Jacobian determinant is $r^2\sin\theta$.
\item[Intermediate 17.] Lagrange multipliers give $(1,1)/\sqrt2$ for the maximum and its negative for the minimum.
\item[Intermediate 23.] For $r\ne0$, direct differentiation yields zero; the origin must be treated separately.
\item[Advanced 6.] Apply Cauchy--Schwarz to $D_{\hat u}f=\nabla f\cdot\hat u$ and identify the equality condition.
\item[Advanced 8.] Symmetry and the multiplier equations imply $|x|=|y|=|z|=1$ at nonzero extrema; inspect signs.
\item[Advanced 11.] Write $F_i=f(r)x_i/r$ and differentiate components before summing.
\item[Advanced 13.] Two differentiations give $\nabla^2e^{i\mathbf k\cdot\mathbf r}=-|\mathbf k|^2e^{i\mathbf k\cdot\mathbf r}$.
\item[Challenge 3.] Vary $q(t)\mapsto q(t)+\epsilon\eta(t)$ with $\eta$ vanishing at both endpoints, differentiate the action with respect to $\epsilon$, and integrate the $\dot\eta$ term by parts.
\item[Challenge 5.] Compute the phase-space divergence of $(\partial H/\partial p_i,-\partial H/\partial q^i)$; equality of mixed partials causes cancellation.
\item[Challenge 8.] Curl annihilates gradients wherever $\chi$ is sufficiently smooth.
\item[Challenge 10.] Apply the divergence theorem to the flux term, then use the arbitrariness of the volume.
\end{description}
